% ─────────────────────────────────────────────────────────────────────────────
% Bibliography entries for the weak memory models section.
% Add these to your main \begin{thebibliography}{...} block.
% ─────────────────────────────────────────────────────────────────────────────

\bibitem[AABN16]{AABN16}
P.~A. Abdulla, M.~F. Atig, A.~Bouajjani, and T.~P. Ngo.
\newblock The benefits of duality in verifying concurrent programs under {TSO}.
\newblock In \emph{Proceedings of CONCUR 2016}, Leibniz International
  Proceedings in Informatics, vol.~59, pp.~5:1--5:15.
  Schloss Dagstuhl, 2016.

\bibitem[AABN18]{AABN_journal}
P.~A. Abdulla, M.~F. Atig, A.~Bouajjani, and T.~P. Ngo.
\newblock A load-buffer semantics for total store ordering.
\newblock \emph{Logical Methods in Computer Science}, 14(1):9, 2018.

\bibitem[AAR20]{AAR20}
P.~A. Abdulla, M.~F. Atig, and R.~Rezvan.
\newblock Parameterized verification under {TSO} is {PSPACE}-complete.
\newblock \emph{Proceedings of the ACM on Programming Languages},
  4(POPL):26:1--26:29, 2020.

\bibitem[AABKS25]{eTSO_param}
P.~A. Abdulla, M.~F. Atig, A.~Bouajjani, K.~N. Kumar, and P.~Saivasan.
\newblock Parametrised verification of {Intel} x86 programs.
\newblock \emph{Proceedings of the ACM on Programming Languages},
  10(POPL), 2026.

\bibitem[AABKS24]{eTSO_kalternation}
P.~A. Abdulla, M.~F. Atig, A.~Bouajjani, K.~N. Kumar, and P.~Saivasan.
\newblock Verification under {Intel}-x86 with persistency.
\newblock \emph{Proceedings of the ACM on Programming Languages},
  8(PLDI):1189--1212, 2024.

\bibitem[ABBM10]{ABBM10}
M.~F. Atig, A.~Bouajjani, S.~Burckhardt, and M.~Musuvathi.
\newblock On the verification problem for weak memory models.
\newblock In \emph{Proceedings of POPL 2010}, pp.~7--18. ACM, 2010.

\bibitem[Lam79]{Lam79}
L.~Lamport.
\newblock How to make a multiprocessor computer that correctly executes
  multiprocess programs.
\newblock \emph{IEEE Transactions on Computers}, C-28(9):690--691, 1979.

\bibitem[RMV22]{eTSO_model}
A.~Raad, L.~Maranget, and V.~Vafeiadis.
\newblock Extending {Intel}-x86 consistency and persistency: Formalising the
  semantics of {Intel}-x86 memory types and non-temporal stores.
\newblock \emph{Proceedings of the ACM on Programming Languages},
  6(POPL):22:1--22:31, 2022.
